% Owner: P1. Keep the final introduction around 0.6 pages.
\section{Introduction}
Passage retrieval selects evidence for question answering from a large corpus.
Dense retrievers encode questions and passages into continuous vectors, enabling
semantic matching beyond exact word overlap. These representations can retrieve
effectively, but their individual dimensions are difficult to interpret, and
storing a dense vector for every passage is expensive.

\citet{park-etal-2025-decoding} introduce sparse autoencoders (SAEs) to represent
dense retrieval embeddings through a small set of active latent concepts.
Concept-Level Sparse Retrieval (CL-SR) searches these representations with an
inverted index. This exposes individual matching concepts and offers a route
to reducing storage. Two objectives remain distinct: reconstructing an embedding
and ranking relevant passages. A low reconstruction error alone does not ensure
that the final concept scores preserve the teacher's retrieval preferences.

Our project asks whether explicitly optimizing ranking, together with controlling
concept usage, improves this quality--efficiency trade-off. We investigate five
stages independently: learned impact scoring (P1), ranking-aware SAE training
(P2), adaptive passage sparsity (P3), concept pruning (P4), and query concept
expansion (P5). Each experiment must use a measured baseline under its own
controlled protocol. Combining results across stages additionally requires
compatible checkpoints, corpus statistics, and evaluation settings.

P1 isolates the scoring rule. A small neural network learns each shared concept's
contribution from dense SimLM scores, while keeping the indexed representations
fixed. Its additive form permits full-corpus concept retrieval. We evaluate
against released and validation-tuned BM25, with held-out queries excluded from
selection. The completed experiment improves top-10 quality, but also exposes
costs in latency and deeper TREC recall. Reporting these costs is necessary to
assess whether a ranking gain is practically useful.
